\documentclass[11pt, twocolumn, landscape, a4paper]{article}

\usepackage[margin=1cm]{geometry}
\usepackage[utf8]{inputenc}
\usepackage[spanish]{babel}
\usepackage{amssymb, amsmath, amsbsy}
\usepackage{tasks}
\usepackage{enumerate}
\usepackage{graphicx}

\usepackage[pdftex]{hyperref}
\hypersetup{pdftitle={Aritmética - Varios},pdfauthor={Luis Carrillo Gutiérrez}}

\fontfamily{cmss}\selectfont
\renewcommand{\familydefault}{\sfdefault}
\pagestyle{empty}

\begin{document}
\subsection*{Aritmética - conjuntos}
\begin{itemize}
\item{Colocar el valor de verdad/falseda en cada proposición para el conjunto:
$A = \{2\,;\,3\,;\,\{1\}\,;\,\{2\,;\,1\}\}$
\begin{enumerate}[I.]
\item{$\varnothing \in A$}
\item{$3 \in A$}
\item{$1 \in A$}
\item{$\{1\} \in A$}
\item{$\{3\} \subset A$}
\item{$\varnothing \subset A$}
\end{enumerate}
\begin{tasks}(5)
\task{FVFFVV}\task{FFVVFF}\task{FFFVVV}\task{FVFVFV}\task{VVFVFV}
\end{tasks}
}
\item{¿Cuántos subconjuntos tiene $A = \{1\,;\,\{1\}\,;\,1\,;\,\varnothing\}$?
\begin{tasks}(5)
\task{4}\task{8}\task{15}\task{16}\task{32}
\end{tasks}
}
\end{itemize}
\end{document}
